% Part: applied-modal-logics
% Chapter: epistemic-logic
% Section: public-announcement-logic-lang

\documentclass[../../../include/open-logic-section]{subfiles}

\begin{document}

\olfileid{aml}{el}{pal}

\olsection{જાહેર જાહેરાતનો તર્ક}

ગતિશીલ જ્ઞાનસંબંધી તર્કપદ્ધતિઓ સમય જતાં અથવા નવી માહિતી મળતાં
કર્તાઓનું જ્ઞાન કેવી રીતે બદલાય છે તે દર્શાવી શકે છે. તેમાંની
ઘણી પદ્ધતિઓ જ્ઞાનમાં થતા ફેરફારોને માહિતીલક્ષી \emph{ઘટનાઓ}
અથવા \emph{અપડેટો} વડે દર્શાવે છે. સૌથી મૂળભૂત અપડેટ
\emph{જાહેર જાહેરાત} છે: કોઈ સૂત્ર સત્યપૂર્વક જાહેર થાય છે
અને બધા કર્તાઓ એકસાથે એ ઘટનાના સાક્ષી બને છે. આ દર્શાવવા
માટે ભાષા નીચે મુજબ વિસ્તરીએ છીએ.

\begin{defn}
$G$ એ કર્તા-સંકેતોનો ગણ હોય. જાહેર જાહેરાતો ધરાવતી બહુ-કર્તા
જ્ઞાનસંબંધી તર્કની મૂળભૂત ભાષામાં નીચેની બાબતો હોય છે:
\begin{enumerate}
  \tagitem{prvFalse}{!!{falsity} માટેનો વિધાનાત્મક અચલાંક~$\lfalse$.}{}
  \tagitem{prvTrue}{!!{truth} માટેનો વિધાનાત્મક અચલાંક~$\ltrue$.}{}
  \item !!{propositional variable}sનો !!{denumerable}s ગણ: $\Obj
    p_0$, $\Obj p_1$, $\Obj p_2$, \dots
  \item વિધાનાત્મક સંયોજકો: \startycommalist
  \iftag{prvNot}{\ycomma $\lnot$ (નિષેધ)}{}%
  \iftag{prvAnd}{\ycomma $\land$ (સંયોજન)}{}%
  \iftag{prvOr}{\ycomma $\lor$ (વિયોજન)}{}%
  \iftag{prvIf}{\ycomma $\lif$ (!!{conditional})}{}%
  \sourcecorrection{OLMD-080}{મૂળના સંયોજકોના કોષ્ટકમાં
  દ્વિશરતી $\liff$ રહી ગયું છે, જ્યારે નીચેની આગમનાત્મક
  વ્યાખ્યામાં તે હાજર છે. અહીં તેને ઉમેર્યું છે.}
  \iftag{prvIff}{\ycomma $\liff$ (!!{biconditional})}{}.
  \item $a \in G$ હોય ત્યારે જ્ઞાન-સંક્રિયક $\Knows_a$.
  \item $!B$ એ !!a{formula} હોય ત્યારે જાહેર-જાહેરાત
  સંક્રિયક $[!B]$.
\end{enumerate}
\end{defn}

જાહેર-જાહેરાત સંક્રિયક બૉક્સ સંક્રિયકની જેમ કાર્ય કરે છે;
તેથી ભાષાની આગમનાત્મક વ્યાખ્યા નીચે મુજબ છે:

\begin{defn}
  જ્ઞાનસંબંધી ભાષાનાં \emph{!!^{formula}s} નીચે મુજબ આગમનાત્મક રીતે
  વ્યાખ્યાયિત થાય છે:
  \begin{enumerate}
  \tagitem{prvFalse}{$\lfalse$ પરમાણુક !!{formula} છે.}{}

  \tagitem{prvTrue}{$\ltrue$ પરમાણુક !!{formula} છે.}{}

  \item દરેક વિધાનાત્મક ચલ $\Obj p_i$ (પરમાણુક) !!{formula} છે.

  \tagitem{prvNot}{જો $!A$ એ !!a{formula} હોય, તો $\lnot !A$ પણ
  !!a{formula} છે.}{}

  \tagitem{prvAnd}{જો $!A$ અને $!B$ એ !!{formula}s હોય, તો $(!A \land
  !B)$ પણ !!a{formula} છે.}{}

  \tagitem{prvOr}{જો $!A$ અને $!B$ એ !!{formula}s હોય, તો $(!A \lor !B)$
  પણ !!a{formula} છે.}{}

  \tagitem{prvIf}{જો $!A$ અને $!B$ એ !!{formula}s હોય, તો $(!A \lif !B)$
  પણ !!a{formula} છે.}{}

  \tagitem{prvIff}{જો $!A$ અને $!B$ એ !!{formula}s હોય, તો $(!A \liff !B)$
  પણ !!a{formula} છે.}{}

  \item જો $!A$ એ !!a{formula} હોય અને $a \in G$ હોય, તો $\Knows_a !A$
  પણ !!a{formula} છે.
  
  \item જો $!A$ અને $!B$ એ !!{formula}s હોય, તો $[!A] !B$ પણ
  !!a{formula} છે.

  \tagitem{limitClause}{આ સિવાય બીજું કશું !!a{formula} નથી.}{}
\end{enumerate}
\end{defn}

!!{formula}~$[!A] !B$ને ``$!A$ની સત્ય જાહેરાત થયા પછી $!B$
સત્ય છે'' એમ વાંચવાનું છે. જાહેર જાહેરાતોના સંદર્ભમાં ક્યારેક
સામાન્ય જ્ઞાન વિશે પણ વાત કરવી ઉપયોગી થશે; તેથી ભાષામાં
સામાન્ય જ્ઞાનનો સંક્રિયક~$\CKnows_G !A$ પણ હોઈ શકે છે.

\end{document}
